% Chapter 3, typeset from lectures/L03: the README is the opener and the closing review, the
% lecture's Appendix A is the chapter's one section of material, and Appendix B its exercises.

\chapter{Interrupts}
\label{c3:ch}

{\large\itshape Letting the hardware say when something has happened, instead of asking it over
and over.\par}

\begin{chapterbox}{In this chapter}
\begin{itemize}
  \item Why polling doesn't scale, and what an interrupt actually is.
  \item External interrupts: INT0/INT1 (\code{SREG}, \code{EICRA}, \code{EIMSK}), edge selection,
    and the \code{ISR}/vector mechanism.
  \item Worked example: toggle an LED via INT0.
  \item Worked example: blink an LED via INT1.
  \item Pin Change Interrupts (PCI): any other pin, at the cost of per-port grouping and
    both-edges triggering.
  \item Worked example: toggle an LED via a PCI interrupt, and the general pattern for any port.
  \item When to reach for INT0/INT1 versus PCI.
\end{itemize}
\end{chapterbox}

Every input the book has read so far has been polled. This chapter is about the alternative: why
interrupts exist, and what problem they solve compared to polling. You will configure an external
interrupt, INT0 or INT1, for a specific edge and write its ISR; configure a Pin Change Interrupt on
an arbitrary pin, including working out the correct \code{PCINTx} number; and see why PCI
interrupts fire on both edges, and how to filter for the one you care about.

Read \secref{c3:sec:interrupts} first: it is the chapter's core material, with its worked
examples. Then work the exercises in \secref{c3:sec:exercises}. The Reflection and Calculation
exercises are answered in Appendix~\ref{app:answers}; the Code exercises have no printed solution,
because the board tells you whether a program does what the exercise asks.

% Section 3.1, from lectures/L03/appendix/a_interrupts.md.

\section{Interrupts}
\label{c3:sec:interrupts}\label{c3:app:a}

\subsection{Why interrupts}
\label{c3:a:1}
\begin{itemize}
  \item Every input we've read so far has been polled: the \code{while} loop checks
    \code{BTN1_PRESSED} every iteration. Fine for one button, but it wastes cycles and doesn't
    scale; imagine polling a dozen inputs, or needing sub-millisecond reaction time on one of them
    while the rest of the program does something else entirely.
  \item An interrupt lets a hardware event (button edge, timer overflow, ADC done, UART byte
    received\ldots) suspend whatever the CPU is doing, run a dedicated handler (an
    \textbf{Interrupt Service Routine}, or \textbf{ISR}), and then resume exactly where it left
    off. It's the embedded equivalent of an event handler, except the ``event loop'' is the CPU
    itself.
  \item Analogy: an MP3 player's main loop plays the current track. Pressing ``next'' doesn't get
    checked between every audio sample; it interrupts playback, the ISR queues the next track, and
    playback resumes.
  \item Because you don't control \emph{when} an ISR fires, keep them short: flip a flag or update
    a small piece of state, and let \code{main()}'s loop do the actual work. We'll follow that
    convention throughout.
\end{itemize}

\subsection{External interrupts: INT0 / INT1}
\label{c3:a:2}
The ATmega328P has two pins wired straight to dedicated interrupt logic: \textbf{INT0} on D2
(\code{PORTD2}) and \textbf{INT1} on D3 (\code{PORTD3}).

Three things to configure:
\begin{enumerate}
  \item \textbf{Global interrupt enable}, the \code{I} bit in the status register \code{SREG}.
    \code{<avr/interrupt.h>} provides \code{sei()}/\code{cli()} for this, thin wrappers around the
    \code{SEI}/\code{CLI} assembly instructions:

\begin{ccode}
sei(); // Enable interrupts globally.
cli(); // Disable interrupts globally.
\end{ccode}
  \item \textbf{Edge/level selection}, per interrupt, in \code{EICRA} (bits
    \code{ISCn1:ISCn0}):

\begin{narrowtable}{@{}ll@{}}
  \thead{\code{ISCn1:ISCn0}} & \thead{Triggers on} \\
  \midrule
  \code{00} & pin is low \\
  \code{01} & any logic change \\
  \code{10} & falling edge (high $\rightarrow$ low) \\
  \code{11} & rising edge (low $\rightarrow$ high) \\
\end{narrowtable}

\begin{ccode}
EICRA = (1U << ISC00) | (1U << ISC01); // Enable INT0 on rising edge.
\end{ccode}
  \item \textbf{Enable the specific interrupt} in \code{EIMSK}:

\begin{ccode}
EIMSK = (1U << INT0);
\end{ccode}
\end{enumerate}

The handler itself is a function named \code{ISR}, parameterized by a \emph{vector} that ties it
to a specific interrupt source, \code{INT0_vect} for INT0, \code{INT1_vect} for INT1. No prototype,
no explicit call site; the vector is how the linker wires it up.

\begin{ccode}
ISR(INT0_vect)
{
    // Runs on the configured edge, then execution resumes exactly where the main program
    // was interrupted.
}
\end{ccode}

\subsection{Example: toggle an LED via INT0}
\label{c3:a:3}
Button on D2 (INT0, rising edge), LED on D9. Every press flips the LED.

\bookimage{lectures/L03/appendix/images/circuit_int0_toggle.png}{A button on D2 (INT0) and an LED
on D9.}{c3:fig:int0}

\begin{ccode}
#include <avr/interrupt.h>
#include <avr/io.h>

/** GPIO pins. */
#define LED1 1U // D9  -> PORTB1.

/** GPIO operations. */
#define LED1_TOGGLE PINB = (1U << LED1) // Toggle LED1.

/**
 * @brief Set up system.
 */
static void setup(void)
{
    // Configure LED1 as output.
    DDRB = (1U << LED1);

    // Enable rising edge interrupt on INT0 (D2) => button pressdown will cause an interrupt.
    EICRA = (1U << ISC00) | (1U << ISC01);
    EIMSK = (1U << INT0);

    // Enable interrupts globally.
    sei();
}

/**
 * @brief Toggle the LED on button event.
 */
ISR(INT0_vect) { LED1_TOGGLE; }

/**
 * @brief Application entry point.
 *
 * @return 0 on termination of the program (should never occur).
 */
int main(void)
{
    setup();
    while (1) {}
    return 0;
}
\end{ccode}

\textbf{Why \code{LED1_TOGGLE} writes to \code{PINB}, not \code{PORTB}.} \code{PINx} registers are
normally read-only: reading \code{PINB} returns the actual logic level currently on each of Port
B's pins, regardless of whether they're configured as inputs or outputs. But writing a \code{1} to a
\code{PINxn} bit is a documented exception: on this chip (and other modern AVR cores), it
\emph{toggles} the corresponding \code{PORTxn} bit instead of doing nothing, independent of that
pin's \code{DDRxn} setting. Only the bits written as \code{1} are affected; any \code{0} bits in
the write leave their \code{PORTx} counterpart untouched.

This isn't a hack, it's a deliberate hardware feature, and it exists because a toggle written the
``obvious'' way, \code{PORTB ^= (1U << LED1);}, is a read-modify-write: the CPU reads \code{PORTB},
flips the bit, then writes the result back. If an interrupt changes another bit in \code{PORTB}
between that read and that write, the interrupt's change gets silently overwritten once the write
completes. Writing to \code{PINB} instead needs no read at all, so there's no window for that race;
the toggle is atomic no matter what else might be touching \code{PORTB}.

INT1 on D3 works identically for configuration; swap \code{ISC00}/\code{ISC01} for
\code{ISC10}/\code{ISC11}, \code{INT0} for \code{INT1} in \code{EIMSK}, and \code{INT0_vect} for
\code{INT1_vect}. \Secref{c3:a:4} walks through a full example.

\subsection{Example: blink an LED via INT1}
\label{c3:a:4}
Button on D3 (INT1, falling edge, i.e.\ on release), LED on D9. Each press/release cycle toggles
between the LED blinking every 100\,ms and staying off.

\bookimage{lectures/L03/appendix/images/circuit_int1_blink.png}{A button on D3 (INT1) and an LED
on D9.}{c3:fig:int1}

\begin{ccode}
#define F_CPU 16000000UL // CPU frequency in Hz.

#include <stdbool.h>
#include <stddef.h>
#include <stdint.h>

#include <avr/interrupt.h>
#include <avr/io.h>
#include <util/delay.h>

/** GPIO pins. */
#define LED1 1U // D9 -> PORTB1.

/** GPIO operations. */
#define LED1_ON PORTB |= (1U << LED1)   // Enable LED1.
#define LED1_OFF PORTB &= ~(1U << LED1) // Disable LED1.

/** Time parameters. */
#define TICK_1MS 1U         // Generate a 1 ms tick in delay_ms().
#define BLINK_SPEED_MS 100U // Blink speed in ms.

/** Blink state (true = blinking, false = off); shared with the ISR, see below. */
static volatile bool blink_state = false;

/** Current blink speed in ms; also shared with the ISR. */
static volatile uint16_t blink_speed_ms = 0U;

/**
 * @brief Set up system.
 */
static void setup(void)
{
    // Configure LED1 as output.
    DDRB = (1U << LED1);

    // Enable falling edge interrupt on INT1 (D3).
    EICRA = (1U << ISC11);
    EIMSK = (1U << INT1);

    // Enable interrupts globally.
    sei();
}

/**
 * @brief Generate delay.
 *
 * @param[in] ms Pointer to the duration in ms.
 *
 * @note A pointer is used so that the duration can be set to 0 on blink state change.
 */
static void delay_ms(const volatile uint16_t* ms)
{
    // Null check included to demonstrate the practice; ms is only ever passed as
    // &blink_speed_ms below, so this never actually triggers. Both this check and the
    // <stddef.h> include it needs (for NULL) can be omitted.
    if (NULL == ms) { return; }

    for (uint16_t i = 0U; i < *ms; ++i)
    {
        _delay_ms(TICK_1MS);
    }
}

/**
 * @brief Blink the LED.
 */
static void led_blink(void)
{
    LED1_ON;
    delay_ms(&blink_speed_ms);
    LED1_OFF;
    delay_ms(&blink_speed_ms);
}

/**
 * @brief Toggle the blink state on button event, update the blink speed accordingly.
 */
ISR(INT1_vect)
{
    blink_state    = !blink_state;
    blink_speed_ms = blink_state ? BLINK_SPEED_MS : 0U;
}

/**
 * @brief Application entry point.
 *
 * @return 0 on termination of the program (should never occur).
 */
int main(void)
{
    setup();

    while (1)
    {
        if (blink_state) { led_blink(); }
        else { LED1_OFF; }
    }
    return 0;
}
\end{ccode}

\textbf{Why \code{blink_state} and \code{blink_speed_ms} are \code{volatile}.} Both are written
inside \code{ISR(INT1_vect)} and read elsewhere: \code{blink_state} in \code{main()}'s loop,
\code{blink_speed_ms} via the pointer \code{delay_ms()} dereferences. Without \code{volatile}, the
compiler is free to assume a variable never changes on its own between two reads in the same
function, and is allowed to cache it in a register once instead of re-reading memory every time; it
has no way to know an interrupt can run between those reads and change the value. Skipping
\code{volatile} here wouldn't necessarily fail right away, either: the loop might work fine at low
optimization levels and only start ignoring button presses once a build enables optimizations
aggressive enough to actually act on that caching assumption, which is exactly the kind of bug
that's miserable to track down later. \code{volatile} tells the compiler to always re-read (or
re-write) the variable from memory, never to assume its value between accesses.
Chapters~\ref{c5:ch} and~\ref{c6:ch} both reuse this same pattern for their own ISR-shared
variables.

\textbf{Why \code{delay_ms()} takes a pointer, not a plain \code{uint16_t}.} \Secref{c1:a:7}
introduced pointers as C's way to let a function reach into the caller's own object instead of
working on a private copy, there, so a function could write an out-parameter (\code{*x = 3;}). This
is the same mechanism, used the other way around: \code{delay_ms(&blink_speed_ms)} doesn't hand the
function \code{blink_speed_ms}'s value, it hands over its \emph{address}, so every \code{*ms} inside
the loop re-reads whatever \code{blink_speed_ms} currently holds, not a value frozen at the moment
\code{delay_ms()} was called.

That distinction matters here specifically because \code{blink_speed_ms} can change out from under
\code{delay_ms()} mid-delay, the ISR overwrites it the instant the button is released. Had
\code{delay_ms()} taken \code{const uint16_t ms} by value instead, it would have received a private
copy at call time and stayed blind to the ISR's update for the rest of the delay. With a pointer,
the very next \code{i < *ms} check after the ISR runs sees the new value, so a press/release is
noticed within about one 1\,ms tick instead of waiting for the whole delay to finish first.

This live-reading trick relies on the same \code{volatile} guarantee explained above: without it,
the compiler could assume \code{*ms} never changes across loop iterations and cache it in a
register instead of re-reading memory each time, silently defeating the exact behavior this
pointer parameter exists to enable.

The blink itself still uses a blocking \code{_delay_ms()} loop, same as the hand-rolled delay of
\secref{c2:a:6}; this chapter only changes \emph{how} the toggle is triggered, not how the blink
timing works. Non-blocking timing is the subject of Chapter~\ref{c5:ch}, on timers.

\subsection{Pin Change Interrupts (PCI) for any other pin}
\label{c3:a:5}
INT0/INT1 only cover two physical pins. Every other pin can still trigger an interrupt via
\textbf{Pin Change Interrupts}, at the cost of some extra bookkeeping:
\begin{itemize}
  \item PCI interrupts are grouped \textbf{by port, not by pin}: all of \code{PORTB} shares one
    vector (\code{PCINT0_vect}), all of \code{PORTC} shares another (\code{PCINT1_vect}), all of
    \code{PORTD} a third (\code{PCINT2_vect}). Enable PCI on two pins in the same port and both
    fire the \emph{same} ISR; the vector alone doesn't tell you which pin changed.
  \item PCI fires on \textbf{both} edges (press and release); there's no edge-select register like
    \code{EICRA} for it. If you only care about one edge, check the pin's actual level inside the
    ISR and ignore the call otherwise.
\end{itemize}
Configuration, per port:

\begin{narrowtable}{@{}llll@{}}
  \thead{Port} & \thead{Enable bit (\code{PCICR})} & \thead{Mask register} &
    \thead{Bit-name formula} \\
  \midrule
  B (D8--D13) & \code{PCIE0} & \code{PCMSK0} & \code{PCINTx = PORTBx} \\
  C (A0--A5) & \code{PCIE1} & \code{PCMSK1} & \code{PCINTx = PORTCx + 8} \\
  D (D0--D7) & \code{PCIE2} & \code{PCMSK2} & \code{PCINTx = PORTDx + 16} \\
\end{narrowtable}

Example: button on D13 (\code{PORTB5}) $\rightarrow$ \code{PCINT5}.

\bookimage{lectures/L03/appendix/images/circuit_pci_toggle.png}{A button on D13 (\code{PCINT5})
and an LED on D9.}{c3:fig:pci}

\begin{ccode}
#include <avr/interrupt.h>
#include <avr/io.h>

/** GPIO pins. */
#define LED1 1U // D9  -> PORTB1.
#define BTN1 5U // D13 -> PORTB5.

/** GPIO operations. */
#define LED1_TOGGLE PINB = (1U << LED1)    // Toggle LED1.
#define BTN1_PRESSED (PINB & (1U << BTN1)) // High if BTN1 is pressed, low otherwise.

/**
 * @brief Set up system.
 */
static void setup(void)
{
    // Configure LED1 as output.
    DDRB = (1U << LED1);

    // Configure BTN1 as input with its internal pull-up enabled.
    PORTB = (1U << BTN1);

    // Enable pin change interrupt on D13 (for BTN1).
    PCICR  = (1U << PCIE0);
    PCMSK0 = (1U << BTN1);

    // Enable interrupts globally.
    sei();
}

/**
 * @brief Toggle the LED on rising button edge.
 */
ISR(PCINT0_vect)
{
    if (BTN1_PRESSED) { LED1_TOGGLE; }
}

/**
 * @brief Application entry point.
 *
 * @return 0 on termination of the program (should never occur).
 */
int main(void)
{
    setup();
    while (1) {}
    return 0;
}
\end{ccode}

\subsection{PCI on any port: the general pattern}
\label{c3:a:6}
The port B example of \secref{c3:a:5} generalizes directly to ports C and D: only the enable bit,
mask register, and vector change, following the same formula table from \secref{c3:a:5}.

\begin{narrowtable}{@{}llll@{}}
  \thead{Port} & \thead{Enable (\code{PCICR})} & \thead{Mask register} & \thead{Vector} \\
  \midrule
  B & \code{PCICR = (1U << PCIE0);} & \code{PCMSK0 = (1U << PCINTx);} & \code{ISR(PCINT0_vect)} \\
  C & \code{PCICR = (1U << PCIE1);} & \code{PCMSK1 = (1U << PCINTx);} & \code{ISR(PCINT1_vect)} \\
  D & \code{PCICR = (1U << PCIE2);} & \code{PCMSK2 = (1U << PCINTx);} & \code{ISR(PCINT2_vect)} \\
\end{narrowtable}

Worked example: a button on D2 (\code{PORTD2}) needs port D's enable bit and mask register, and its
ISR is tied to \code{PCINT2_vect}:

\begin{ccode}
#define BTN2 2U // D2 -> PORTD2.

static void setup(void)
{
    // ...
    PCICR = (1U << PCIE2);  // Enable PCI on port D.
    PCMSK2 = (1U << BTN2);  // Specifically D2 / PORTD2.
    sei();                  // Enable interrupts globally.
}

ISR(PCINT2_vect) // Shared by every enabled PCI pin on port D.
{
    // ...
}
\end{ccode}

\code{PCMSKn}'s bit position always matches the pin's own bit position within its port
(\code{DDRx}/\code{PORTx}/\code{PINx}); the \code{PCINTx = PORTDx + 16}-style formula of
\secref{c3:a:5} only matters for computing the \emph{global} \code{PCINTx} number (useful for
looking a pin up in the datasheet), not for what bit to actually set in \code{PCMSKn}. So
\code{PCMSK2 = (1U << BTN2)} above sets exactly the same bit as avr-libc's named
\code{PCMSK2 = (1U << PCINT18)} would; defining your own pin-bit macro, the same convention
\secref{c3:a:5} already used for \code{BTN1}, skips the global-number lookup entirely. Same idea
for a button on D7 (\code{PORTD7}): \code{PCMSK2 = (1U << 7)} sets the same bit as
\code{(1U << PCINT23)}, no formula needed.

The same button on A0 (\code{PORTC0}) instead needs \code{PCIE1} / \code{PCMSK1}, and an
\code{ISR(PCINT1_vect)}.

\subsection{Choosing between INT0/INT1 and PCI}
\label{c3:a:7}
\begin{itemize}
  \item Need a specific edge (rising \emph{or} falling, not both) and don't mind being limited to
    D2/D3? Use INT0/INT1; it's simpler, and the edge filtering happens in hardware.
  \item Need any other pin, or don't care which edge? Use PCI, and filter edges yourself by reading
    the pin level in the ISR, as above.
  \item Multiple PCI sources on the same port share an ISR, and if you have more than one, the
    handler needs to work out which pin actually changed (e.g.\ by keeping a ``previous state''
    byte and XOR-ing it against the current \code{PINx} read to see which bit flipped).
\end{itemize}

% Section 3.2, from lectures/L03/README.md: the goals, the evaluation questions, and the next
% lecture.

\section{Review}
\label{c3:sec:review}

\needspace{6\baselineskip}
\subsection*{What you should be able to do now}
After this chapter, you should:
\begin{itemize}
  \item Explain, in your own words, why interrupts exist and what problem they solve compared to
    polling.
  \item Be able to configure an external interrupt (INT0 or INT1) for a specific edge and write its
    ISR.
  \item Be able to configure a Pin Change Interrupt on an arbitrary pin, including working out the
    correct \code{PCINTx} number.
  \item Understand why PCI interrupts fire on both edges, and how to filter for the one you care
    about.
\end{itemize}

\needspace{6\baselineskip}
\subsection*{Questions to test yourself}
You should be able to answer:
\begin{enumerate}
  \item What does \code{EICRA} configure that \code{EIMSK} doesn't, and vice versa?
  \item Why do multiple Pin Change Interrupts on the same port share a single ISR, and what does
    that mean for how you write it?
  \item A PCI-based button handler toggles an LED on both press \emph{and} release. What line is
    missing, and why?
\end{enumerate}

\needspace{6\baselineskip}
\subsection*{Looking ahead}
Chapter~\ref{c4:ch} is about analog-to-digital conversion:
\begin{itemize}
  \item Analog-to-digital conversion (ADC): reading a continuous signal instead of a simple
    high/low.
  \item A worked example combining what you already know (PCI, bitwise register I/O) with the ADC.
\end{itemize}

% Section 3.3, from lectures/L03/appendix/b_exercises.md.

\section{Exercises}
\label{c3:sec:exercises}\label{c3:app:b}

\begin{aside}
\textbf{How to work these.} The Reflection and Calculation exercises are answered on paper, from
the chapter. The Code exercises are programs for the board: build and flash them the way
\secref{c2:app:c} (Microchip Studio) or \secref{c2:app:d} (Linux) describes, or run them in the
simulator of Appendix~\ref{sim:ch} if you have no board yet.

\textbf{How to check your work.} The Reflection and Calculation exercises are answered in
Appendix~\ref{app:answers}. The Code exercises have no printed solution: the board tells you
whether a program does what the exercise describes. Try each exercise before you read its answer.
\end{aside}

% ------------------------------------------------------------------------------------------
\exerciseset{Concepts}
\label{c3:set:1}

\reflexercise{Polling versus interrupts}
\label{c3:ex:1}
In your own words (2--3 sentences), explain why polling \code{BTN1_PRESSED} in a \code{while} loop
scales poorly once you have several inputs to watch, and how interrupts solve that.

\reflexercise{\code{EICRA} and \code{EIMSK}}
\label{c3:ex:2}
What does \code{EICRA} configure that \code{EIMSK} doesn't? Could you enable INT0 without touching
\code{EICRA} at all; what would happen if you did?

% ------------------------------------------------------------------------------------------
\exerciseset{INT0/INT1}
\label{c3:set:2}

\codeexercise{Moving the button to INT1}
\label{c3:ex:3}
Rewire the INT0 example of \secref{c3:a:3} (button on D2) to use \textbf{INT1} on D3 instead.
Write out the three configuration lines that change.

\reflexercise{The wrong edge}
\label{c3:ex:4}
A colleague configures \code{EICRA} for a falling edge instead of rising, but leaves the rest of
the INT0 example unchanged, and complains the LED toggles when the button is \emph{released}, not
pressed. Explain what happened.

% ------------------------------------------------------------------------------------------
\exerciseset{Pin Change Interrupts}
\label{c3:set:3}

\calcexercise{An interrupt on D4}
\label{c3:ex:5}
A project needs an interrupt on D4. Work out, from the formula table in \secref{c3:a:5}, which
\code{PCINT} number that is, which port group it belongs to, and which three registers/bits need to
be touched to enable it.

\reflexercise{Two buttons, one vector}
\label{c3:ex:6}
Two buttons are wired to D11 and D12 (both \code{PORTB}). Both are configured via PCI and share
\code{PCINT0_vect}. Sketch how a single ISR could tell which button changed, using a ``previous
state'' byte compared against the current \code{PINB} read.

% ------------------------------------------------------------------------------------------
\exerciseset{Program}
\label{c3:set:4}

\codeexercise{A PCI blink toggle}
\label{c3:ex:7}
Wire up an LED on D6 and a button on D4. Using a Pin Change Interrupt (not INT0/INT1) on the
button, write a complete program, from scratch, that:
\begin{itemize}
  \item Starts with the LED off and blinking disabled.
  \item Each button press toggles between two states: blinking (LED on/off every 200\,ms) and
    disabled (LED held off).
  \item The ISR only flags that the button was pressed, filtering the release edge the same way as
    the PCI example of \secref{c3:a:5}; the actual state toggle and the blink timing both happen in
    \code{main()}.
\end{itemize}
You'll need a blocking delay for the blink timing itself (\code{_delay_ms()}, as in
Chapter~\ref{c2:ch}); non-blocking delays are the subject of Chapter~\ref{c5:ch}, on timers, not
of this one. This chapter's own \secref{c3:a:4} (blink an LED via INT1) is the closest worked
example to lean on.

\codeexercise{An interrupt-driven state machine}
\label{c3:ex:8}
Build a 5-state state machine, driven entirely by interrupts, using every kind covered in this
chapter at once.

Wire up three LEDs, \code{LED1} on D8, \code{LED2} on D9, and \code{LED3} on D10, plus five
buttons. Declare the five states as an \code{enum} (\code{STATE_OFF}, \code{STATE_SLOW},
\code{STATE_MEDIUM}, \code{STATE_FAST}, \code{STATE_ON}), and wire each button as follows,
advancing/rewinding with wraparound at both ends of the range:
\begin{itemize}
  \item \code{BTN1} on D2 (\textbf{INT0}, falling edge) is a reset: jump straight to
    \code{STATE_OFF}, ignoring the normal advance/rewind stepping.
  \item \code{BTN2} on D3 (\textbf{INT1}, rising edge) rewinds the state by two steps.
  \item \code{BTN3} on D11 (\textbf{PCI}) rewinds the state by one step.
  \item \code{BTN4} on D12 (\textbf{PCI}) advances the state by one step.
  \item \code{BTN5} on D13 (\textbf{PCI}) advances the state by two steps.
\end{itemize}
The three LEDs' behavior depends on the current state:
\begin{itemize}
  \item \code{STATE_OFF}: all three LEDs held off.
  \item \code{STATE_SLOW} / \code{STATE_MEDIUM} / \code{STATE_FAST}: the three LEDs chase in a loop
    (\code{LED1} $\rightarrow$ \code{LED2} $\rightarrow$ \code{LED3} $\rightarrow$ \code{LED1}
    \ldots), lighting exactly one at a time, holding each for 500\,ms / 250\,ms / 50\,ms
    respectively.
  \item \code{STATE_ON}: all three LEDs held on.
\end{itemize}
As with every other example in this chapter, keep each ISR minimal: flag which button fired (and,
for the three PCI buttons sharing one vector, work out \emph{which} pin changed, as in
Exercise~\ref{c3:ex:6}), and let \code{main()} do the actual work of updating \code{state} and
driving the LEDs.

\bookimage{lectures/L03/appendix/images/circuit_state_machine_exercise.jpg}{Three LEDs on D8--D10
and five buttons, on D2, D3 and D11--D13.}{c3:fig:statemachine}

